\section{问题分析与总体路线}
\subsection{关键权衡}
合并强依赖算子可以减少边界 COPY、重复读取和任务释放等待，但过度聚合会压缩多核并行空间，并可能让核内启发式在 L1/UB 容量压力下插入更多换入换出。过度切分则使共享 DDR 竞争、同步及启动间隔增加。因此，通信字节数、任务数或各核算子数都不是总完成时间的充分代理。若 Cube 与 Vector 工作量不均，按算子个数平衡也会让其中一条 Pipe 长期空闲。论文以合法方案的官方 Makespan 为最终准则，将成本估计只用于有限候选的生成和排序。

三问共同采用“结构候选—依赖调度—官方评价—有界改进”的框架。先识别计算依赖和共享输入，构造完整连通分量、主导分量有界拆分等不同粒度的合法初解；再按后续关键工作量和最早可行完成时间安排核心与空隙；最后用少量局部划分、迁移和顺序调整寻找更短的总时间。场景 B 的候选必须重新按同核驻留、跨核 COPY 释放和容量规则评价；共享 L2 的候选还需经过 FIFO 事件仿真。算法保留已知可行方案，在限定评价机会内比较，不进行分组、核心和排列的组合穷举。

\subsection{三种运行语义的关系}
表\ref{tab:scenario} 概括后两问相对前一问的改变。相同的输出字段不意味着相同的代价函数：尤其在场景 B 中，一个核心内的多个子图仍有先后次序，但其通信可利用驻留，跨核释放以具体数据搬运事件为准。问题三缓存只改变合资格读取的服务路径，不能绕过算子依赖、同步和四条 Pipe 的占用。
\begin{table}[!htbp]\centering
\caption{三问运行语义与增量机制}\label{tab:scenario}
\begin{tabularx}{\linewidth}{p{0.12\linewidth}XXX}
\toprule
内容 & 问题一：场景 A & 问题二：场景 B & 问题三：B 加共享 L2\\
\midrule
Task 与驻留 & 一子图一 Task；边界不驻留 & 同核子图同属一 Task；可驻留 & 保持 B 的 Task 与私有存储\\
跨核数据 & 边界经共享 DDR；按 A 释放 & COPY 完成后另有同步释放 & 同 B，不因命中取消同步\\
读取服务 & DDR & DDR & COPY\_IN 可命中共享 FIFO L2\\
主要决策 & 划分、映射、核上序 & 同左，重估驻留和生命周期 & 同左，重估缓存事件与双带宽池\\
\bottomrule
\end{tabularx}
\end{table}
